% !TEX program = lualatex
\documentclass[a4paper,10pt,twocolumn]{ltjsarticle}
\usepackage{amsmath,amssymb,booktabs,geometry,graphicx,hyperref,url}
\geometry{top=18mm,bottom=20mm,left=16mm,right=16mm,columnsep=7mm}
\renewcommand{\dbltopfraction}{0.95}
\renewcommand{\topfraction}{0.9}
\renewcommand{\textfraction}{0.05}
\hypersetup{colorlinks=true,linkcolor=blue,citecolor=blue,urlcolor=blue}
\graphicspath{{notes/figures/supervisor_comparison_20261007/}}
\newcommand{\NRMSE}{\mathrm{NRMSE}}
\newcommand{\HV}{\mathrm{HV}}
\newcommand{\clip}{\mathrm{clip}}
\title{問題設定3案の予備実験と\\状態観測あり・なしBO比較の報告}
\author{}
\date{実験日：2026年10月1日\quad 資料作成日：2026年10月7日}
\begin{document}
\maketitle
\begin{abstract}
本研究は，多目的遺伝的プログラミング（GP）の探索状態に応じて，交叉率，突然変異率，更新周期をベイズ最適化（BO）で調整し，収束性と多様性維持の両立を目指す．本資料では，今後の比較問題を選ぶためのUBall5D・Airfoil・Concrete各3回の予備実験と，セミナー最終条件のFriedman-Iで状態入力の効果を調べた100回の比較実験を報告する．予備実験ではAirfoilに後半の改善余地が見られ，UBall5Dでは定数式への停滞，Concreteでは手法による停滞の違いが見られた．状態入力の比較では，平均最終HVの優位性は確認されなかった一方，提案法のHV標準偏差は約26\%小さかった．問題の継続選定と状態情報の効果の検証を分け，今後の検証方針を整理する．
\end{abstract}

\section{研究目的と今回の報告範囲}
GPは，入力と出力の関係を表す数式を，演算子と変数からなる木として探索する手法である．本研究では，予測誤差と式の複雑さを同時に小さくする多目的GPを用いる．例えば，小さく読みやすい式と，大きいが誤差の小さい式の両方を候補として残し，精度と複雑さのトレードオフを得る．

提案法の目的は，複数の世代状態に基づいて操作率と更新タイミングを調整し，収束性と多様性維持の改善を図ることである．GPを制御対象，BOを制御器とみなし，交叉率\(p_c\)，突然変異率\(p_m\)，次回調整までの世代数\(k\)を決定する．\(k\)世代の実行後に成果を評価し，次の設定を選ぶ閉ループである．

今回の報告は，次の二つの問いに対応する．第一は，\textbf{ウォームアップ後にも探索が進み，制御の効果を調べられる問題はどれか}である．UBall5D・Airfoil・Concreteの3案を，既存のセミナー条件へ移植して各3回実行した．第二は，\textbf{状態をBOへ入力すること自体に追加的な効果があるか}である．Friedman-Iの最終セミナー条件で，提案法と非文脈BOを100回ずつ比較した．前者は問題選定のための予備観察，後者は状態入力を除く比較であり，結果の根拠の強さを区別する．

\section{共通の手法と評価指標}
\subsection{多目的GPとBO制御}
全手法で，同じNSGA-II型の世代更新を用いる．親集団と子集団を併合し，他の解に支配されない解を優先して残し，同じ非劣ランクでは混雑距離を用いて目的空間の広がりを保つ．探索中に得られた非劣解を保存する外部集合をアーカイブと呼ぶ．アーカイブは成果評価に用い，多様性や平均木サイズは現在集団から観測する．

提案法がBOへ渡す状態は，世代進行率，アーカイブHV，直近HV改善量，構造多様性，平均木サイズ，停滞長の6成分である．現在状態\(x_\ell\)に応じた区間報酬を
\begin{equation}
 r=f(x_\ell,p_c,p_m,k)
\end{equation}
として学習する．更新周期は\(k\in\{1,3,5\}\)とし，各\(k\)に1本のガウス過程回帰モデルを持つ．各モデルで操作率候補の期待改善量（EI）を評価し，次の制御入力を選ぶ．

実験で実際に用いた報酬は
\begin{equation}
 r_\ell=0.55h_\ell+0.35a_\ell-0.05c_\ell-0.25b_\ell
\end{equation}
である．\(h_\ell\)は世代あたりHV改善量を0.01で正規化し\([-1,1]\)へ制限した値，\(a_\ell\)は区間平均多様性と目標との一致度，\(c_\ell\)は短い更新周期へのコスト，\(b_\ell\)は多様性が0.10を下回るときの罰則である．目標多様性は進行度に応じて0.60から0.20へ下げる．したがって，多様性項は\textbf{多様性が高いほど常に高報酬になる項ではない}．停滞・木の肥大化への罰則の重みは現行条件で0である．

\subsection{指標の意味と評価予算}
訓練データでの目的は
\begin{equation}
 f_1(T)=\NRMSE_{\mathrm{train}}(T),\qquad f_2(T)=|T|
\end{equation}
の2目的最小化である．\(T\)は式木，\(|T|\)はノード数である．NRMSEは二乗平均平方根誤差を目標値の標準偏差で割った値で，小さいほど良い．訓練平均を常に出す定数式の訓練NRMSEは1となるため，1付近の値は単純な平均予測から十分改善していない可能性を示す．

HVは，参照点に対して解集合が支配する面積で，大きいほど精度と複雑さのトレードオフが良い．今回のHV計算では，NRMSEを\([0,5]\)，木サイズを\([1,100]\)に基づいて正規化し，参照点を\((1,1)\)とした．\textbf{1ノード・NRMSE=1の定数式でもHVは0.8になり得るため，HV=0.8を80\%の予測精度と読んではいけない}．100ノードは正規化の基準であり，木のハード上限ではない．

集団多様性\(D\)は，木の部分構造の共有度に基づく個体間距離の平均で，0に近いほど似た個体が多い．高い\(D\)自体は良い予測を保証しないため，HV・誤差・式の内容と併せて読む．以下の平均と標準偏差は実行間で計算し，標準偏差の分母は実行回数\(n\)（\texttt{ddof=0}）とする．標準偏差は平均値の信頼区間とは異なる．

両実験とも24個体，ウォームアップ18世代，本評価60世代，合計78世代である．BOの初期区間は\(k=1,1,3,3,5,5\)で，合計18世代を消費する．固定率GPも78世代実行し，18世代を同じ比較境界とするが，BOのウォームアップは行わない．1実行の評価回数は子集団分1872回，初期集団24回を含めると1896回である．検証・テストの検索後評価はこの予算に含めない．

\section{比較用問題設定3案の予備実験}
\subsection{目的・選定方針・共通設定}
Friedman-Iでは初期から18世代までに改善が集中していたため，BOによる制御を評価する残りの探索余地が小さい可能性があった．そこで，真の式が分かる合成問題，停滞機序を調べやすい実データ，異なる分野の実データという3種類を選んだ．問題の選定理由は，\textbf{提案法が勝つことではなく，再現性と探索挙動の違いを観察できること}である．Whiteらのベンチマーク提案\cite{white2013}と，Liuらの精度・サイズを目的とする多目的GPの研究\cite{liu2022}を根拠とした．

\begin{table}[tb]
\centering\small
\caption{3問題の予備実験で比較した手法}
\label{tab:methods}
\begin{tabular}{@{}lcc@{}}\toprule
手法 & \(p_c\) & \(p_m\)\\\midrule
標準固定率 & 0.80 & 0.05\\
高突然変異固定率 & 0.70 & 0.20\\
高交叉固定率 & 0.90 & 0.05\\
提案法 & BOで調整 & BOで調整\\\bottomrule
\end{tabular}
\end{table}

表\ref{tab:methods}の4手法を，各問題について乱数シード0，1，2で実行した（3問題\(\times\)4手法\(\times\)3回＝36実行）．各問題で，同じデータと初期集団を全手法へ与えた．データ生成・分割の乱数シードは20260527で固定し，探索用の乱数と分離した．3回は\textbf{同じデータ分割上の探索反復}であり，3通りのデータ分割ではない．

関数集合は四則演算，保護付き除算，\(\sin\)，\(\cos\)で，定数生成範囲は\([-2,2]\)，初期木深さ5，突然変異で生成する部分木の深さ2，親選択のトーナメントサイズ3とした．アーカイブ上限は200である．保護付き除算は分母の絶対値が\(10^{-6}\)以下なら分子を返し，予測値を\(\pm10^6\)へ制限する．BO・報酬・生存選択は既存のセミナー実装を維持した．原論文の大集団，関数の追加，係数の事後最適化を再現した実験ではない．

\subsection{UBall5D：5変数の関係を組み合わせる合成問題}
UBall5D（Vladislavleva-4）は
\begin{equation}
 y=\frac{10}{5+\sum_{i=1}^{5}(x_i-3)^2}
\end{equation}
を目標とする．5入力すべてを分母の中で組み合わせる必要があり，真の式が既知なので誤差と構造の両面を解釈できる．WhiteらのTable 5・6.1節に従い，訓練1024点を各入力\(U[0.05,6.05]\)，テスト5000点を\(U[-0.25,6.35]\)から生成した\cite{white2013}．独立した検証1024点を訓練と同じ範囲から追加した点は本研究独自の設定である．ノイズは加えない．

選定理由は，変数の組合せと外挿を含み，単純な関数より構造探索の継続を観察できる可能性があることである．ただし，難しい問題でも，少数個体では単に改善不能な停滞へ入る場合がある．「長く改善しない」ことを「長期探索に適した問題」と取り違えない必要がある．

\subsection{Airfoil：翼の騒音を予測する実データ}
Airfoil Self-Noiseは，周波数，迎角，翼弦長，流速，吸込み側の排除厚さの5入力から音圧レベルを予測する1503例の公開データである\cite{airfoil}．LiuらはAirfoilを例に，小さい式の過剰複製によって大きい式が生まれにくくなる探索挙動を示している（3.1節・図1，6.1.2節・図3）\cite{liu2022}．多様性，平均木サイズ，停滞を状態に含める本研究と対応するため，制御の機序を調べる候補とした．

今回は訓練901例，検証225例，テスト377例に固定分割し，入力・出力を訓練データだけの平均・標準偏差で標準化した．原研究は生存選択も変更しているため，その改善が本研究の操作率制御だけで再現できるとはいえない．原研究と集団サイズ・関数集合も異なり，元論文の曲線から本実験の収束時期を予測しない．

\subsection{Concrete：配合と材齢から強度を予測する実データ}
Concrete Compressive Strengthは，セメント，高炉スラグ，フライアッシュ，水，減水剤，粗骨材，細骨材，材齢の8入力から圧縮強度を予測する1030例の公開データである\cite{concrete}．Liuらの多目的GP実験にも用いられている（Table 2・3）\cite{liu2022}．Airfoilと異なる物理現象で検証でき，入力数と標本数の面でも扱いやすいことを理由に選んだ．

今回は訓練614例，検証156例，テスト260例とした．同一入力の重複34行を削除せず，同じ入力の行を同一分割にまとめて情報の漏れを防いだ．そのため行数比率は厳密な60：15：25ではない．標準化はAirfoilと同様に訓練データだけから計算した．8入力であること自体は，長期に改善する根拠ではない．

\subsection{結果の概観}
表\ref{tab:pilot}に最終HV，集団多様性，18世代後のHV増加量を示す．図\ref{fig:pilot}は，これらに加えて訓練誤差がどの世代で改善したかを確認する図である．提案法の平均最終HVが4手法中で最良だったのはAirfoilである．UBall5Dでは標準固定率，Concreteでは高突然変異固定率が最良だった．ただし，各3回の平均であり，統計的優位を示す比較ではない．

\begin{table*}[tb]
\centering\small
\caption{3問題の予備実験結果（各手法3回，平均\(\pm\)標準偏差）}
\label{tab:pilot}
\input{notes/figures/supervisor_comparison_20261007/pilot_metrics_rows.tex}
\end{table*}

\begin{figure*}[tb]
\centering\includegraphics[width=\textwidth]{pilot_dynamics.png}
\caption{3問題の世代推移．各行はUBall5D，Airfoil，Concrete，各列はアーカイブHV，集団多様性，訓練誤差．線は3回の平均，帯は標準偏差，破線は18世代の比較境界を表す．訓練誤差は各世代の訓練誤差最小の式について記録する．多様性は構造の異なりを表し，誤差の良さとは別の指標である．}
\label{fig:pilot}
\end{figure*}

\subsection{各問題で分かったことと考察}
\paragraph{UBall5D：高いHVだけでは学習成功といえない．}
提案法の平均最終HVは0.798186だが，訓練誤差最小の最終式は3回とも入力変数を使わない定数式で，訓練NRMSEは平均1.000449だった．18世代後HV増加は平均0.018752だが，最も増加した実行でも最終式は定数である．したがって，HV改善を真の5変数関係の学習と解釈できない．標準固定率では1回だけ12ノードの非定数式が得られたが，利用変数は1個にとどまった．

真の式を現在の関数集合と定数の組合せで100ノード以内に構成し，訓練・テストともNRMSEが\(10^{-12}\)未満になる確認は済んでいる．表現不能というより，現行の少数集団と誤差・サイズの選択がその構造へ到達していない可能性がある．現条件のまま長期主問題に採用せず，まず固定率だけで集団サイズや小木への偏りを調べる．

\paragraph{Airfoil：後半の改善を検証する有力候補．}
提案法の18世代後HV増加は，乱数シード順に0.027427，0.016051，0.002104で3回とも正であり，うち2回は60世代以降にも改善した．高突然変異固定率との最終HV差は平均\(+0.018369\)で，3回すべて上回った．標準固定率と高交叉固定率にも後半の改善があり，提案法だけが改善する問題ではない．

一方，標準固定率の18世代後HV増加平均0.017507は，提案法の0.015194より大きい．また，表\ref{tab:error}では提案法の訓練・検証誤差は最小だが，テスト誤差は標準固定率の方が小さい．したがって，現段階では「継続探索を観察できる候補」として優先し，汎化性能や提案法の全面的優位が示されたとは扱わない．

\paragraph{Concrete：早期停滞を診断する副候補．}
提案法の2回は，標準化したセメント量\(x_1\)のみを使う6ノード式
\[\sin(\sin(\sin(\sin(\sin(x_1)))))\]
で停止し，18世代後のHV増加は0だった．残り1回の増加は0.003089で，平均増加は0.001030である．一方，高突然変異固定率の2回では18世代後に0.033050，0.003279の改善があり，有限予算内の改善90\%到達世代は74，71だった．

したがって，問題自体が18世代で必ず解けるわけではなく，探索条件によって後半改善の有無が変わる．提案法は最終HV・訓練誤差・検証誤差で高突然変異固定率を上回っていない．長期主問題としての優先度は下げるが，小さい式への停滞と制御の限界を調べる副実験として残す．原因を突然変異率や報酬だけに特定する証拠はまだない．

\begin{table*}[tb]
\centering\small
\caption{訓練誤差最小の最終式のNRMSE（各3回，平均\(\pm\)標準偏差）}
\label{tab:error}
\input{notes/figures/supervisor_comparison_20261007/pilot_error_rows.tex}
\end{table*}

\subsection{問題選定と結果解釈の注意点}
表\ref{tab:error}は，探索終了後に訓練誤差最小の1本を未見データへ適用した診断値であり，テスト上のパレートフロントやテストHVではない．テストが最良の式を選び直してはいない．各分割のNRMSEはその分割の目標値の標準偏差で計算し，検索用の\([0,5]\)制限は加えずに報告する．UBall5Dのテストは範囲を広げた外挿であり，実データの固定分割とは難しさが異なる．

早期改善の集中度を確認する補助指標として，18世代までの改善割合\(q_W\)と，78世代までの改善の90\%に達する世代\(t_{90}\)を用いた．提案法の中央値はUBall5Dで96.7\%・13世代，Airfoilで32.8\%・63世代，Concreteで100\%・13世代である．これらは\textbf{有限予算内の改善の配分}であり，真の収束や318世代での挙動の証明ではない．

本番問題は提案法の勝敗で決定せず，未使用の探索乱数で固定率のみを再評価し，後半の改善，非定数式の獲得，状態による操作率の有効性の変化を確認してから選ぶ．現時点での次の検証優先度はAirfoil，Concrete，UBall5Dである．3回・単一分割では，実行乱数やデータ分割に対する一般性を評価できない．

表\ref{tab:pilot}の多様性についても，提案法が全問題・全固定率を上回る傾向は見られない．例えばAirfoilの提案法は平均0.6053で，標準固定率0.6262や高交叉固定率0.6416より低い．UBall5Dでは比較的高い多様性でも最良式は定数であり，構造の違いが有用な変数関係の獲得を保証しないことが分かる．また，問題ごとに出力の分布や前処理が異なるため，異なる問題間の絶対HVを難度の順位として比較しない．

\section{セミナー最終条件での状態観測あり・なしBO比較}
\subsection{目的と公平な比較の設定}
以前の固定率比較では，提案法は標準固定率GPを上回った．しかし，この差には報酬を用いた操作率探索，ウォームアップ，\(k\)調整などが含まれ，状態情報の寄与は切り分けられない．そこで，BOへ入力する世代状態だけを取り除く比較を行った．

対象はFriedman-Iの
\begin{equation}
\begin{aligned}
 y={}&10\sin(\pi x_1x_2)+20(x_3-0.5)^2\\
     &+10x_4+5x_5,\quad x_i\sim U(0,1)
\end{aligned}
\end{equation}
で，訓練200点，テスト200点を用意した．今回の主評価は訓練目的に基づくHVと集団多様性で，テスト誤差は主比較へ用いていない．24個体・78世代・乱数シード0～99の100回を各条件で実行した．

提案法は状態6成分と操作率2成分の計8次元を代理モデルへ入力する．\textbf{非文脈BO（本資料でいう観測なしBO）}は操作率の2次元のみを入力し，学習対象を
\begin{equation}
 r=f(p_c,p_m,k)
\end{equation}
とする．非文脈BOも実際に得た報酬を観測して学習する．状態の計測は報酬と評価のため両条件で継続するので，\textbf{状態入力の利用効果の比較であり，計測を省く費用削減の比較ではない}．

各\(k\)のモデル，EI比較，候補128点／\(k\)，初期設計点各2回，報酬，率の範囲と変化幅制約，平滑化係数0.60を共通とした．カーネルは定数倍のARD付きMat\'ern（\(\nu=2.5\)）と白色雑音の和で，入力次元に合わせた長さ尺度の数のみ8から2へ変更した．残り世代数による\(k\)の実行可能性判定も共通である．

全100回で最初の6区間の制御記録が完全一致し，同じ18世代後集団から比較できていることを確認した．提案法の全世代履歴も，6月3日の既存セミナー最終結果と完全一致した．よって，今回の提案法側は別条件の再実験結果ではなく，既存最終条件の再現である．

\subsection{結果：平均性能の差とばらつきの差を区別}
表\ref{tab:ablation}と図\ref{fig:ablation-progress}に結果を示す．平均最終HVは提案法0.802461，非文脈BO0.802270で，差は\(+0.000191\)（相対約\(+0.024\%\)）だった．同じ乱数シードを対応させた平均差の95\%信頼区間は\([-0.005478,\,+0.005859]\)で0を含む．対応あり\(t\)検定は\(p=0.947\)，Wilcoxon符号付順位検定は\(p=0.408\)で，\textbf{状態入力による平均最終HVの優位性は確認できなかった}．差が検出されないことは，同等性が証明されたこととは異なる．

\begin{table*}[tb]
\centering\small
\caption{状態入力の有無による比較（各100回，平均\(\pm\)標準偏差）}
\label{tab:ablation}
\input{notes/figures/supervisor_comparison_20261007/ablation_metrics_rows.tex}
\end{table*}

\begin{figure*}[tb]
\centering\includegraphics[width=\textwidth]{ablation_dynamics.png}
\caption{Friedman-IにおけるアーカイブHVと集団多様性の推移．線は100回の平均，帯は標準偏差．破線の18世代までは両条件の制御記録が一致している．平均曲線は近い一方で，提案法の実行間のばらつきは比較的小さい．}
\label{fig:ablation-progress}
\end{figure*}

本評価中の平均HVは，世代18～78のHV曲線を台形積分して60で割った時間方向の平均水準である．提案法0.787353，非文脈BO0.785468で，差の95\%信頼区間は\([-0.003917,\,+0.007687]\)，対応あり\(t\)検定は\(p=0.521\)だった．平均的な収束過程の改善も明確ではない．最終多様性の平均も0.640513と0.638855で近く，多様性改善を強く主張する結果ではない．

一方，最終HV標準偏差は0.022454と0.030345で，提案法が約26.0\%小さい．同じ乱数シードの対を単位に20,000回再標本化したブートストラップでは，標準偏差差（提案法－非文脈BO）の95\%区間は\([-0.016234,\,-0.000504]\)だった．\textbf{今回の100回ではばらつき低減を支持する}が，他問題への一般化や低下の原因は未検証である．本評価中の平均HVの標準偏差も約38.2\%小さかった．

\begin{figure*}[tb]
\centering\includegraphics[width=\textwidth]{ablation_distribution.png}
\caption{最終HVの分布と同一乱数シードの対応比較．左の点は全100回，箱は第1・第3四分位，箱内の線は中央値，ひげは四分位範囲の1.5倍以内の範囲を表す．右は1点が同じ乱数シードの対で，対角線より上なら提案法のHVが高い．}
\label{fig:ablation-distribution}
\end{figure*}

図\ref{fig:ablation-distribution}の対応比較では，提案法は43勝・5引分・52敗で，差の中央値は\(-0.000902\)だった．全実行で一様に改善したのではない．非文脈BOの最小HVは0.652827で提案法の0.751448より低い一方，最大HVは非文脈BOの0.869018が提案法の0.853518より高い．5\%分位は0.773619と0.772076，10\%分位は0.777284と0.777234で近い．ばらつき差には少数の大きな失敗と高成績側の広がりが影響しており，「すべての低成績実行を救済した」とは解釈しない．

\subsection{制御挙動とその解釈}
図\ref{fig:k}の本評価中の\(k=1\)の選択割合は，提案法56.7\%，非文脈BO48.2\%だった．\(k=5\)は20.0\%と27.0\%で，提案法の方が短い区間を多く選んだ．これらは制御区間数の割合であり，各\(k\)で費やした世代数の割合ではない．本評価中の制御更新回数は平均26.48回と23.30回で，提案法が3.18回多かった（差の95\%信頼区間\([+0.7995,+5.5605]\)）．

\begin{figure}[tb]
\centering\includegraphics[width=\linewidth]{ablation_k_selection.png}
\caption{本評価中の更新周期の選択割合．100回の制御区間を合算し，各条件の区間数で割った値である．}
\label{fig:k}
\end{figure}

短周期化で状態変化に対応する機会が増え，結果を安定させた可能性はある．ただし，今回の比較では操作率と\(k\)の両方が変わるので，\(k\)だけが安定化の原因とは断定できない．また，更新回数が増えており，計算時間の優位や制御回数の削減は示していない．各実行の区間平均操作率をさらに100回で平均すると，\(p_c\)は0.7136と0.6954，\(p_m\)は0.1403と0.1410だった．突然変異率の平均は近いが，平均値だけで状態への応答を説明することはできない．

\subsection{状態情報の効果に関する考察}
第一に，両BOの平均最終HVは約0.802で，過去の標準固定率0.782991より高いが，高突然変異固定率0.808315には届かない．したがって，\textbf{固定率を上回ったことと，状態情報の利用が有効だったことは分けて結論づける必要がある}．今回の追加比較により，状態入力の上乗せ効果は平均性能よりもばらつき低減として現れた可能性が明らかになった．

第二に，共通の18世代で，提案法の初期から最終までの平均HV改善の約82.2\%を得ている．本評価区間で残る探索余地が小さく，状態ごとの違いが有効になる場面が少なかった可能性がある．ただし，ウォームアップ後の改善量も両条件で近いため，最終HVの評価だけが差を隠したとはいえない．

第三に，提案法は8次元，非文脈BOは2次元の代理モデル入力を持つが，本評価中の報酬観測は平均26.48回と23.30回で，さらに\(k\)別モデルに分かれる．状態依存性を表現する利点と，少量データでは学習しにくい不利が共存している可能性がある．これは説明仮説であり，学習不足そのものを測定した結果ではない．

第四に，報酬には多様性目標との一致と更新コストも含まれるため，報酬の最大化が最終HV最大化と一致するとは限らない．本評価の全制御区間を等重みで平均した報酬は，提案法0.047159，非文脈BO0.059945だったが，区間長や探索時期が異なるため直接の性能順位には使わない．今回の比較では報酬項を変更していない．

\section{結論と次に進める検証}
3問題の予備実験では，Airfoilで固定率を含めた後半改善が見られ，次の継続検証の有力候補となった．Concreteは早期停滞と操作条件の違いを調べる副候補，UBall5Dは定数式への縮退を解消できるか確認する候補として保持する．この優先順位は今回の条件での作業判断であり，本番採用や統計的優劣の確定ではない．

Friedman-Iの100回比較では，\textbf{状態入力による平均最終HV・平均収束過程・多様性の明確な改善は示されなかった一方，最終HVのばらつきは低減した}．現時点では，提案法の目的を全面的に達成したと結論づけず，「状態依存制御が探索結果の安定化へ寄与する可能性を示した」と整理する．

次の検証は三つに絞る．第一に，Airfoilを中心に未使用の探索乱数で固定率のみの長期予備実験を行い，改善が続く条件を，提案法の勝敗を見ずに確定する．第二に，同じ問題・予算・ウォームアップで状態入力あり／なしを比較し，平均性能，多様性，ばらつきの再現性を調べる．第三に，\(k\)を揃えた比較によって状態入力と短周期化の寄与を切り分け，必要に応じて報酬項の個別比較へ進む．3問題の今回使用済みのテスト結果を本番の未使用評価とは扱わず，探索反復とデータ分割の不確実性も分けて検証する．

\section*{結果と文献の追跡情報}
今回の実験は2026年10月1日に実行済みで，本資料作成時に再実行していない．数値は保存済み世代履歴・制御履歴・集計JSONから転記し，図表は同じ記録から再作成した．3問題の詳細は\path{notes/thesis_three_problem_pilot_report.md}，状態入力比較の詳細は\path{notes/seminar_context_ablation_report.md}にある．元の結果フォルダはそれぞれ\path{outputs/thesis_three_problem_pilot/seminar_conditions_seed3_20261001/}，\path{outputs/seminar_context_ablation/20261001/}である．

参考文献[1]のTable 5・6.1節をUBall5Dの式・標本仕様と問題選定の注意に，[2]の3.1節・図1および6.1.2節・図3をAirfoilの停滞機序の選定理由に，Table 2・3をConcreteの利用実績に参照した．[3][4]はデータ内容と配布元の出典である．分割比率，標準化，集団24，78世代，各3回／100回は本研究の設定であり，原論文の推奨値や完全再現条件ではない．本文確認済み文献を再利用し，問題設定用文献フォルダの原資料を保持した．

\begin{thebibliography}{4}
\footnotesize
\bibitem{white2013}
D. R. White et al., ``Better GP Benchmarks: Community Survey Results and Proposals,''
\emph{Genetic Programming and Evolvable Machines}, 14(1), pp. 3--29, 2013.
DOI: \href{https://doi.org/10.1007/s10710-012-9177-2}{10.1007/s10710-012-9177-2}.
（研究ノート[R29]）
\bibitem{liu2022}
D. Liu, M. Virgolin, T. Alderliesten, and P. A. N. Bosman,
``Evolvability Degeneration in Multi-Objective Genetic Programming for Symbolic Regression,''
\emph{Proc. GECCO}, pp. 973--981, 2022.
DOI: \href{https://doi.org/10.1145/3512290.3528787}{10.1145/3512290.3528787}.
（研究ノート[R30]）
\bibitem{airfoil}
T. Brooks, D. Pope, and M. Marcolini, \emph{Airfoil Self-Noise}, UCI Machine Learning Repository, 1989.
DOI: \href{https://doi.org/10.24432/C5VW2C}{10.24432/C5VW2C}.
（保存データ取得日：2026年10月1日）
\bibitem{concrete}
I. Yeh, \emph{Concrete Compressive Strength}, UCI Machine Learning Repository, 1998.
DOI: \href{https://doi.org/10.24432/C5PK67}{10.24432/C5PK67}.
（保存データ取得日：2026年10月1日）
\end{thebibliography}
\end{document}
